\documentclass[UTF8,12pt]{ctexart}
\usepackage[a4paper,margin=24mm]{geometry}
\usepackage{amsmath,amssymb,bm,graphicx,adjustbox}
\newcommand{\fitmath}[1]{\begin{adjustbox}{max width=\linewidth}$\displaystyle #1$\end{adjustbox}}
\setlength{\parindent}{0pt}
\setlength{\parskip}{.7em}
\sloppy
\allowdisplaybreaks
\begin{document}
\section*{1990 年考研数学三 · 真题}
\subsection*{1990年全国硕士研究生招生考试试题}

\subsection*{（试卷IV）}

\subsection*{一、填空题（本题共5小题，每小题3分，满分15分）}

（1）极限 \(\displaystyle \lim_{n\to\infty}\left(\sqrt{n+3\sqrt n}-\sqrt{n-\sqrt n}\right)=\)\_\_\_\_\_\_。


（2）设函数 \(\displaystyle f(x)\) 有连续的导函数，\(\displaystyle f(0)=0\) 且 \(\displaystyle f\prime(0)=b\)，若函数


\[\fitmath{\displaystyle F(x)=\begin{cases}\dfrac{\displaystyle f(x)+a\sin x}{\displaystyle x},&x\ne0,\\A,&x=0\end{cases}}\]


在 \(\displaystyle x=0\) 处连续，则常数 \(\displaystyle A=\)\_\_\_\_\_\_。


（3）曲线 \(\displaystyle y=x^2\) 与直线 \(\displaystyle y=x+2\) 所围成的平面图形的面积为\_\_\_\_\_\_。


（4）若线性方程组


\[\fitmath{\displaystyle \begin{cases}x_1+x_2=-a_1,\\x_2+x_3=a_2,\\x_3+x_4=-a_3,\\x_4+x_1=a_4\end{cases}}\]


有解，则常数 \(\displaystyle a_1,\allowbreak a_2,\allowbreak a_3,\allowbreak a_4\) 应满足条件\_\_\_\_\_\_。


（5）一射手对同一目标独立地进行4次射击，若至少命中一次的概率为 \(\displaystyle \dfrac{\displaystyle 80}{\displaystyle 81}\)，则该射手的命中率为\_\_\_\_\_\_。


\subsection*{二、选择题（本题共5小题，每小题3分，满分15分）}

（1）设函数 \(\displaystyle f(x)=x\tan x\,e^{\sin x}\)，则 \(\displaystyle f(x)\) 是（　　）


（A）偶函数。 （B）无界函数。 （C）周期函数。 （D）单调函数。


（2）设函数 \(\displaystyle f(x)\) 对任意的 \(\displaystyle x\) 均满足等式 \(\displaystyle f(1+x)=af(x)\)，且有 \(\displaystyle f\prime(0)=b\)，其中 \(\displaystyle a,\allowbreak b\) 为非零常数，则（　　）


（A）\(\displaystyle f(x)\) 在 \(\displaystyle x=1\) 处不可导。 （B）\(\displaystyle f(x)\) 在 \(\displaystyle x=1\) 处可导，且 \(\displaystyle f\prime(1)=a\)。 （C）\(\displaystyle f(x)\) 在 \(\displaystyle x=1\) 处可导，且 \(\displaystyle f\prime(1)=b\)。 （D）\(\displaystyle f(x)\) 在 \(\displaystyle x=1\) 处可导，且 \(\displaystyle f\prime(1)=ab\)。


（3）向量组 \(\displaystyle \boldsymbol\alpha_1,\allowbreak \boldsymbol\alpha_2,\allowbreak \ldots,\allowbreak \boldsymbol\alpha_s\) 线性无关的充分条件是（　　）


（A）\(\displaystyle \boldsymbol\alpha_1,\allowbreak \ldots,\allowbreak \boldsymbol\alpha_s\) 均不为零向量。 （B）其中任意两个向量的分量不成比例。 （C）其中任意一个向量均不能由其余 \(\displaystyle s-1\) 个向量线性表示。 （D）其中有一部分向量线性无关。


（4）设 \(\displaystyle A,\allowbreak B\) 为两随机事件，且 \(\displaystyle B\subset A\)，则下列式子正确的是（　　）


（A）\(\displaystyle P(A+B)=P(A)\)。 （B）\(\displaystyle P(AB)=P(A)\)。 （C）\(\displaystyle P(B\mid A)=P(B)\)。 （D）\(\displaystyle P(B-A)=P(B)-P(A)\)。


（5）设随机变量 \(\displaystyle X\) 和 \(\displaystyle Y\) 相互独立，其概率分布为


\[\fitmath{\displaystyle \begin{array}{c|cc}m&-1&1\\\hline P\{X=m\}&\dfrac{\displaystyle 1}{\displaystyle 2}&\dfrac{\displaystyle 1}{\displaystyle 2}\end{array}\qquad\begin{array}{c|cc}m&-1&1\\\hline P\{Y=m\}&\dfrac{\displaystyle 1}{\displaystyle 2}&\dfrac{\displaystyle 1}{\displaystyle 2}\end{array}}\]


则下列式子正确的是（　　） （A）\(\displaystyle X=Y\)。 （B）\(\displaystyle P\{X=Y\}=0\)。 （C）\(\displaystyle P\{X=Y\}=\dfrac{\displaystyle 1}{\displaystyle 2}\)。 （D）\(\displaystyle P\{X=Y\}=1\)。


\subsection*{三、计算题（本题共4小题，每小题5分，满分20分）}

（1）求函数 \(\displaystyle I(x)=\displaystyle \int_e^x\dfrac{\displaystyle \ln t}{\displaystyle t^2-2t+1}\,dt\) 在区间 \(\displaystyle [e,\allowbreak e^2]\) 上的最大值。


（2）计算二重积分 \(\displaystyle \iint_D xe^{-y^2}\,dx\,dy\)，其中 \(\displaystyle D\) 是曲线 \(\displaystyle y=4x^2\) 和 \(\displaystyle y=9x^2\) 在第一象限所围成的区域。


（3）求级数 \(\displaystyle \sum_{n=1}^{\infty}\dfrac{\displaystyle (x-3)^n}{\displaystyle n^2}\) 的收敛域。


（4）求微分方程 \(\displaystyle y\prime+y\cos x=(\ln x)e^{-\sin x}\) 的通解。


\subsection*{四、（本题满分9分）}

某公司可通过电台及报纸两种方式做销售某种商品的广告，根据统计资料，销售收入 \(\displaystyle R\)（万元）与电台广告费用 \(\displaystyle x_1\)（万元）及报纸广告费用 \(\displaystyle x_2\)（万元）之间的关系有如下经验公式：


\[\fitmath{\displaystyle R=15+14x_1+32x_2-8x_1x_2-2x_1^2-10x_2^2.}\]


（1）在广告费用不限的情况下，求最优广告策略；（2）若提供的广告费用为1.5万元，求相应的最优广告策略。


\subsection*{五、（本题满分6分）}

设 \(\displaystyle f(x)\) 在闭区间 \(\displaystyle [0,\allowbreak c]\) 上连续，其导数 \(\displaystyle f\prime(x)\) 在开区间 \(\displaystyle (0,\allowbreak c)\) 内存在且单调减少，\(\displaystyle f(0)=0\)。试应用拉格朗日中值定理证明不等式 \(\displaystyle f(a+b)\le f(a)+f(b)\)，其中常数 \(\displaystyle a,\allowbreak b\) 满足条件 \(\displaystyle 0\le a\le b\le a+b\le c\)。


\subsection*{六、（本题满分8分）}

已知线性方程组


\[\fitmath{\displaystyle \begin{cases}x_1+x_2+x_3+x_4+x_5=a,\\3x_1+2x_2+x_3+x_4-3x_5=0,\\x_2+2x_3+2x_4+6x_5=b,\\5x_1+4x_2+3x_3+3x_4-x_5=2.\end{cases}}\]


（1）\(\displaystyle a,\allowbreak b\) 为何值时，方程组有解？（2）方程组有解时，求出方程组的导出组的一个基础解系；（3）方程组有解时，求出方程组的全部解。


\subsection*{七、（本题满分5分）}

已知对于 \(\displaystyle n\) 阶方阵 \(\displaystyle A\)，存在自然数 \(\displaystyle k\)，使得 \(\displaystyle A^k=O\)。试证明矩阵 \(\displaystyle E-A\) 可逆，并写出其逆矩阵的表达式（\(\displaystyle E\) 为 \(\displaystyle n\) 阶单位阵）。


\subsection*{八、（本题满分6分）}

设 \(\displaystyle A\) 为 \(\displaystyle n\) 阶矩阵，\(\displaystyle \lambda_1\) 和 \(\displaystyle \lambda_2\) 是 \(\displaystyle A\) 的两个不同的特征值，\(\displaystyle \boldsymbol x_1,\allowbreak \boldsymbol x_2\) 是分别属于 \(\displaystyle \lambda_1\) 和 \(\displaystyle \lambda_2\) 的特征向量。试证明 \(\displaystyle \boldsymbol x_1+\boldsymbol x_2\) 不是 \(\displaystyle A\) 的特征向量。


\subsection*{九、（本题满分4分）}

从0，1，2，\(\displaystyle \ldots\)，9等十个数字中任意选出三个不同的数字，试求下列事件的概率：\(\displaystyle A_1=\{\text{三个数字中不含0和5}\}\)；\(\displaystyle A_2=\{\text{三个数字中不含0或5}\}\)；\(\displaystyle A_3=\{\text{三个数字中含0但不含5}\}\)。


\subsection*{十、（本题满分5分）}

一电子仪器由两个部件构成，以 \(\displaystyle X\) 和 \(\displaystyle Y\) 分别表示两个部件的寿命（单位：千小时），已知 \(\displaystyle X\) 和 \(\displaystyle Y\) 的联合分布函数为


\[\fitmath{\displaystyle F(x,y)=\begin{cases}1-e^{-0.5x}-e^{-0.5y}+e^{-0.5(x+y)},&x\ge0,y\ge0,\\0,&\text{其他}.\end{cases}}\]


（1）问 \(\displaystyle X\) 和 \(\displaystyle Y\) 是否独立？（2）求两个部件的寿命都超过100小时的概率 \(\displaystyle a\)。


\subsection*{十一、（本题满分7分）}

某地抽样调查结果表明，考生的外语成绩（百分制）近似服从正态分布，平均成绩为72分，96分以上的占考生总数的2.3\%，试求考生的外语成绩在60分至84分之间的概率。附表（表中 \(\displaystyle \Phi(x)\) 是标准正态分布函数）


\[\fitmath{\displaystyle \begin{array}{c|ccccccc}x&0&0.5&1.0&1.5&2.0&2.5&3.0\\\hline\Phi(x)&0.500&0.692&0.841&0.933&0.977&0.994&0.999\end{array}}\]


\subsection*{（试卷V）}

\subsection*{一、填空题（本题共5小题，每小题3分，满分15分）}

（1）【同试卷IV 第一、（1）题】 （2）【同试卷IV 第一、（2）题】 （3）【同试卷IV 第一、（3）题】 （4）【同试卷IV 第一、（4）题】


（5）已知随机变量 \(\displaystyle X\sim N(-3,\allowbreak 1),\allowbreak \ Y\sim N(2,\allowbreak 1)\)，且 \(\displaystyle X,\allowbreak Y\) 相互独立，设随机变量 \(\displaystyle Z=X-2Y+7\)，则 \(\displaystyle Z\sim\)\_\_\_\_\_\_。


\subsection*{二、选择题（本题共5小题，每小题3分，满分15分）}

（1）【同试卷IV 第二、（1）题】 （2）【同试卷IV 第二、（2）题】 （3）【同试卷IV 第二、（3）题】


（4）设 \(\displaystyle A\) 为 \(\displaystyle n\) 阶可逆矩阵，\(\displaystyle A^*\) 是 \(\displaystyle A\) 的伴随矩阵，则（　　）


（A）\(\displaystyle |A^*|=|A|^{n-1}\)。 （B）\(\displaystyle |A^*|=|A|\)。 （C）\(\displaystyle |A^*|=|A|^n\)。 （D）\(\displaystyle |A^*|=|A^{-1}|\)。


（5）已知随机变量 \(\displaystyle X\) 服从二项分布，且 \(\displaystyle E(X)=2.4,\allowbreak D(X)=1.44\)，则二项分布的参数 \(\displaystyle n,\allowbreak p\) 的值为（　　） （A）\(\displaystyle n=4,\allowbreak p=0.6\)。 （B）\(\displaystyle n=6,\allowbreak p=0.4\)。 （C）\(\displaystyle n=8,\allowbreak p=0.3\)。 （D）\(\displaystyle n=24,\allowbreak p=0.1\)。


\subsection*{三、（本题共4小题，每小题5分，满分20分）}

（1）求极限 \(\displaystyle \lim_{x\to\infty}\dfrac{\displaystyle 1}{\displaystyle x}\displaystyle \int_0^x(1+t^2)e^{t^2-x^2}\,dt\)。


（2）求不定积分 \(\displaystyle \int\dfrac{\displaystyle x\cos^4(x/2)}{\displaystyle \sin^3x}\,dx\)。


（3）设 \(\displaystyle x^2+z^2=y\varphi(z/y)\)，其中 \(\displaystyle \varphi\) 为可微函数，求 \(\displaystyle \dfrac{\displaystyle \partial z}{\displaystyle \partial y}\)。


（4）【同试卷IV 第三、（2）题】


\subsection*{四、（本题满分9分）【同试卷IV 第四题】}

\subsection*{五、（本题满分6分）}

证明不等式 \(\displaystyle 1+x\ln(x+\sqrt{1+x^2})\ge\sqrt{1+x^2}\)，\(\displaystyle -\infty<x<+\infty\)。


\subsection*{六、（本题满分4分）}

设 \(\displaystyle A\) 为 \(\displaystyle 10\times10\) 矩阵


\[\fitmath{\displaystyle A=\begin{pmatrix}0&1&0&\cdots&0&0\\0&0&1&\cdots&0&0\\\vdots&\vdots&\vdots&\ddots&\vdots&\vdots\\0&0&0&\cdots&0&1\\10^{10}&0&0&\cdots&0&0\end{pmatrix},}\]


计算行列式 \(\displaystyle |A-\lambda E|\)，其中 \(\displaystyle E\) 为10阶单位矩阵，\(\displaystyle \lambda\) 为常数。


\subsection*{七、（本题满分5分）}

设方阵 \(\displaystyle A\) 满足条件 \(\displaystyle A^{\mathrm T}A=E\)，其中 \(\displaystyle A^{\mathrm T}\) 是 \(\displaystyle A\) 的转置矩阵，\(\displaystyle E\) 为单位阵。试证明 \(\displaystyle A\) 的实特征向量所对应的特征值的绝对值等于1。


\subsection*{八、（本题满分8分）【同试卷IV 第六题】}

\subsection*{九、（本题满分5分）【同试卷IV 第九题】}

\subsection*{十、（本题满分6分）}

甲乙两人独立地各进行两次射击，假设甲的命中率为0.2，乙的为0.5，以 \(\displaystyle X\) 和 \(\displaystyle Y\) 分别表示甲和乙的命中次数，试求 \(\displaystyle X\) 和 \(\displaystyle Y\) 的联合概率分布。


\subsection*{十一、（本题满分7分）【同试卷IV 第十一题】}
\end{document}
